\subsubsection{Operator Functions}

Operator functions turn scalar functions into operator-level constructions. Polynomial functional calculus, reduction modulo characteristic polynomials, polynomial spectral mapping, matrix exponentials, and square roots all use the spectral or nilpotent structure developed earlier to compute new operators from old ones.

\begin{theorem} \label{thm:polynomial-functional-calculus-diagonalizable-matrix} \hlt{Polynomial Functional Calculus for Diagonalizable Matrices}\\
Let $A\in\operatorname{Mat}(n,n,\mathbb{F})$ be diagonalizable and write $A=S\Lambda S^{-1}$, where $\Lambda=\operatorname{diag}(\lambda_1,\ldots,\lambda_n)$.\\
If $p\in\mathbb{F}[z]$, then $p(A)=S\operatorname{diag}(p(\lambda_1),\ldots,p(\lambda_n))S^{-1}$.
\begin{proof}
For every $j\geq0$, $A^j=S\Lambda^jS^{-1}$.\\
If $p(z)=\sum_j c_jz^j$, then $p(A)=\sum\limits_j c_jA^j =S\left(\sum\limits_j c_j\Lambda^j\right)S^{-1} =S p(\Lambda) S^{-1}$.\\
Since $\Lambda$ is diagonal, $p(\Lambda)=\operatorname{diag}(p(\lambda_1),\ldots,p(\lambda_n))$.\\
\hlt{Structural Commentary}: Diagonalization reduces $p(A)$ to scalar evaluation on eigenvalues.\\
\hlt{Key Insight}: Compute functions of $A$ by applying the function to the eigenvalues.
\end{proof}
\end{theorem}

\begin{corollary} \label{cor:commutant-distinct-eigenvalues-polynomials} \hlt{Simple Spectrum Commutants Are Polynomials}\\
If $A\in\operatorname{Mat}(n,n,\mathbb{C})$ has $n$ distinct eigenvalues, then $\operatorname{Comm}(A)=\{p(A):p\in\mathbb{C}[z]\}$, $\dim\operatorname{Comm}(A)=n$.
\begin{proof}
Write $A=S\operatorname{diag}(\lambda_1,\ldots,\lambda_n)S^{-1}$ with distinct $\lambda_i$. Theorem \ref{thm:commutant-of-diagonalizable-matrix} says every $B\in\operatorname{Comm}(A)$ has the form $S\operatorname{diag}(\xi_1,\ldots,\xi_n)S^{-1}$. Choose the Lagrange interpolation polynomial $p(z)=\sum\limits_{i=1}^{n}\xi_i\prod\limits_{j\neq i}\frac{z-\lambda_j}{\lambda_i-\lambda_j}$.\\
Then $p(\lambda_i)=\xi_i$ for every $i$, so Theorem \ref{thm:polynomial-functional-calculus-diagonalizable-matrix} gives $B=p(A)$.\\
The dimension statement is the case $n_1=\cdots=n_n=1$ of Theorem \ref{thm:commutant-of-diagonalizable-matrix}.\\
\hlt{Structural Commentary}: Simple spectrum leaves one scalar choice per eigendirection.\\
\hlt{Key Insight}: Interpolation converts arbitrary diagonal data into a polynomial in $A$.
\end{proof}
\end{corollary}

\begin{corollary} \label{cor:polynomial-reduction-by-characteristic-polynomial} \hlt{Reduction Modulo the Characteristic Polynomial}\\
Let $A\in\operatorname{Mat}(n,n,\mathbb{C})$ and let $f\in\mathbb{C}[z]$. If $f=p_Aq+r$ with $\deg r<n$, then $f(A)=r(A)$.
\begin{proof}
Evaluate division identity at $A$: $f(A)=p_A(A)q(A)+r(A)=0q(A)+r(A)=r(A)$.\\
Theorem \ref{thm:cayley-hamilton-matrix} gives $p_A(A)=0$.\\
\hlt{Structural Commentary}: Cayley--Hamilton makes polynomial expressions in $A$ reducible modulo $p_A$.\\
\hlt{Key Insight}: Every high power of $A$ can be rewritten using powers below $n$.
\end{proof}
\end{corollary}

\begin{corollary} \label{cor:powers-of-matrix-span-dimension-at-most-n} \hlt{Powers Span at Most $n$ Dimensions}\\
Let $A\in\operatorname{Mat}(n,n,\mathbb{C})$. Then $\operatorname{span}(I_n,A,A^2,\ldots)=\operatorname{span}(I_n,A,\ldots,A^{n-1})$, so this has dimension at most $n$.
\begin{proof}
For each $k\geq0$, Corollary \ref{cor:polynomial-reduction-by-characteristic-polynomial} applied to $f(z)=z^{n+k}$ gives a polynomial $r_k$ of degree at most $n-1$ such that $A^{n+k}=r_k(A)$. Thus every power of $A$ lies in $\operatorname{span}(I_n,A,\ldots,A^{n-1})$.\\
\hlt{Structural Commentary}: Polynomial reduction collapses the infinite power list to its first $n$ terms.\\
\hlt{Key Insight}: Cayley--Hamilton gives the recurrence relation for powers of $A$.
\end{proof}
\end{corollary}

\begin{corollary} \label{cor:inverse-from-characteristic-polynomial} \hlt{Inverse from the Characteristic Polynomial}\\
Let $A\in\operatorname{Mat}(n,n,\mathbb{C})$ have $p_A(z)=z^n+c_{n-1}z^{n-1}+\cdots+c_1z+c_0$.\\
If $A$ is invertible, then $A^{-1}=-c_0^{-1}(A^{n-1}+c_{n-1}A^{n-2}+\cdots+c_2A+c_1I_n)$.
\begin{proof}
Since $A$ is invertible, $0\notin\operatorname{spec}(A)$, so $c_0=p_A(0)=(-1)^n\det A\neq0$.\\
Cayley--Hamilton gives $A^n+c_{n-1}A^{n-1}+\cdots+c_1A+c_0I_n=0$.\\
Move $c_0I_n$ to the other side and multiply by $-c_0^{-1}A^{-1}$.\\
\hlt{Structural Commentary}: The characteristic equation contains an identity term exactly when $A$ is invertible.\\
\hlt{Key Insight}: Solve Cayley--Hamilton for $A^{-1}$ after factoring out one $A$.
\end{proof}
\end{corollary}

\begin{proposition} \label{prop:eigenvalues-of-polynomial-in-matrix} \hlt{Eigenvalues of a Polynomial in a Matrix}\\
Let $A\in\operatorname{Mat}(n,n,\mathbb{C})$, let $p\in\mathbb{C}[z]$, and let $\lambda_1,\ldots,\lambda_n$ be the eigenvalues of $A$, counted with algebraic multiplicity. Then the eigenvalues of $p(A)$ are $p(\lambda_1),\ldots,p(\lambda_n)$, counted with algebraic multiplicity.
\begin{proof}
Choose a Schur form $U^*AU=T$ with diagonal $\lambda_1,\ldots,\lambda_n$. Then $p(A)=Up(T)U^*$, and $p(T)$ is upper triangular with diagonal $p(\lambda_1),\ldots,p(\lambda_n)$. Hence $p(A)$ has precisely these eigenvalues.\\
\hlt{Structural Commentary}: Schur form reduces the claim to the diagonal.\\
\hlt{Key Insight}: A polynomial preserves triangular form and applies $p$ to each diagonal entry.
\end{proof}
\end{proposition}

\begin{corollary} \label{cor:polynomial-in-matrix-invertible-spectrum-test} \hlt{Polynomial Invertibility Test}\\
Let $A\in\operatorname{Mat}(n,n,\mathbb{C})$ and $p\in\mathbb{C}[z]$. If $p(\lambda)\neq0$ for every $\lambda\in\operatorname{spec}(A)$, then $p(A)$ is invertible.
\begin{proof}
By Proposition \ref{prop:eigenvalues-of-polynomial-in-matrix}, $0\notin\operatorname{spec}(p(A))$. Corollary \ref{cor:matrix-invertible-iff-zero-not-spectrum} gives invertibility.\\
\hlt{Structural Commentary}: Invertibility of $p(A)$ is checked on spectral values of $A$.\\
\hlt{Key Insight}: Avoiding zero on the spectrum prevents $p(A)$ from having a kernel.
\end{proof}
\end{corollary}

\begin{proposition} \label{prop:exponential-of-jordan-block} \hlt{Exponential of a Jordan Block}\\
For $k\geq1$, $\lambda\in\mathbb{C}$, and $t\in\mathbb{C}$, $e^{tJ_k(\lambda)} =e^{\lambda t}\sum\limits_{r=0}^{k-1}\frac{t^r}{r!}J_k(0)^r$.
\begin{proof}
Since $J_k(\lambda)=\lambda I_k+J_k(0)$ and the two summands commute, $e^{tJ_k(\lambda)}=e^{\lambda tI_k}e^{tJ_k(0)}$.\\
The first factor is $e^{\lambda t}I_k$, and the second exponential truncates because $J_k(0)^k=0$.\\
\hlt{Structural Commentary}: The scalar part supplies $e^{\lambda t}$; the nilpotent part supplies a finite polynomial in $t$.\\
\hlt{Key Insight}: Nilpotence turns the exponential series into a finite sum.
\end{proof}
\end{proposition}

\begin{method} \label{meth:solve-linear-ode-by-jordan-form} \hlt{Solving $\mathbf{x}'(t)=A\mathbf{x}(t)$ by Jordan Form}\\
To solve $\mathbf{x}'(t)=A\mathbf{x}(t)$ with $\mathbf{x}(0)=\mathbf{x}_0$:
\begin{enumerate}[label=(\roman*)]
\setlength{\itemsep}{0pt}
\item Put $A=SJS^{-1}$, where $J=\bigoplus_{i=1}^rJ_{k_i}(\lambda_i)$.
\item Set $\mathbf{y}(t)=S^{-1}\mathbf{x}(t)$ and $\mathbf{y}_0=S^{-1}\mathbf{x}_0$. Then $\mathbf{y}'(t)=J\mathbf{y}(t)$ and $\mathbf{x}(t)=S\mathbf{y}(t)$.
\item Split $\mathbf{y}=(\mathbf{y}_1,\ldots,\mathbf{y}_r)$ according to the Jordan blocks. Then $\mathbf{y}_i'(t)=J_{k_i}(\lambda_i)\mathbf{y}_i(t)$.
\item For each block, use Proposition \ref{prop:exponential-of-jordan-block}: $\mathbf{y}_i(t) =e^{tJ_{k_i}(\lambda_i)}\mathbf{y}_i(0) =e^{\lambda_i t}\sum\limits_{q=0}^{k_i-1}\frac{t^q}{q!}J_{k_i}(0)^q\mathbf{y}_i(0)$.
\item Reassemble $\mathbf{y}(t)$ and transform back: $\mathbf{x}(t)=S\mathbf{y}(t)$.
\end{enumerate}
Thus a size-$k$ block with eigenvalue $\lambda$ contributes entries of the form $e^{\lambda t}p(t)$, where $\deg p<k$.
\end{method}

\begin{lemma} \label{lem:identity-plus-nilpotent-square-root} \hlt{Identity Plus Nilpotent Has a Square Root}\\
Let $\mathcal{V}$ be a complex vector space and let $N\in\LL(\mathcal{V})$ be nilpotent. Then $I_{\mathcal{V}}+N$ has a square root.
\begin{proof}
Choose $m\geq1$ such that $N^m=0$. If $m=1$, then $N=0$ and $I_{\mathcal{V}}$ is a square root.\\
Otherwise seek $R=I_{\mathcal{V}}+a_1N+\cdots+a_{m-1}N^{m-1}$. With $a_0=1$, choose $a_1,\ldots,a_{m-1}$ recursively so the coefficients of $R^2$ match those of $I_{\mathcal{V}}+N$: $2a_1=1$, and $2a_k+\sum\limits_{r=1}^{k-1}a_ra_{k-r}=0\quad (2\leq k\leq m-1)$\\
Terms of degree $\geq m$ vanish because $N^m=0$, so $R^2=I_{\mathcal{V}}+N$.\\
\hlt{Structural Commentary}: Nilpotence reduces the square-root series to a finite polynomial.\\
\hlt{Key Insight}: Match coefficients recursively.
\end{proof}
\end{lemma}

\begin{figure}[H]
  \centering
  \begin{tikzpicture}[
  >=Stealth,
  every node/.style={font=\small},
  box/.style={draw=gray!60, rounded corners=3pt, fill=gray!5, inner sep=5pt, align=center},
  arrow/.style={->, line width=0.85pt, blue!70!black}
]
  \node[box] (nil) at (0,0) {$N^m=0$};
  \node[box] (poly) at (3.5,0) {$R=I+a_1N+\cdots+a_{m-1}N^{m-1}$};
  \node[box] (sq) at (7.0,0) {$R^2=I+N$};

  \draw[arrow] (nil) -- (poly);
  \draw[arrow] (poly) -- (sq);
\end{tikzpicture}

  \caption{Nilpotence truncates the square-root series; coefficients are chosen so $R^2=I+N$.}
  \label{fig:nilpotent-square-root-polynomial}
\end{figure}

\begin{theorem} \label{thm:complex-invertible-operator-has-square-root} \hlt{Invertible Complex Operators Have Square Roots}\\
Let $\mathcal{V}$ be finite-dimensional over $\mathbb{C}$ and let $T\in\LL(\mathcal{V})$ be invertible. Then $T$ has a square root.
\begin{proof}
Let $\mathcal{U}_j=\mathcal{G}_{\lambda_j}(T)$ for the distinct eigenvalues of $T$. Primary decomposition gives $\mathcal{V}=\mathcal{U}_1\oplus\cdots\oplus\mathcal{U}_m$, and invertibility gives $\lambda_j\neq0$. Put $N_j=(T-\lambda_jI_{\mathcal{V}})|_{\mathcal{U}_j}$.\\
Then $T|_{\mathcal{U}_j} =\lambda_j\left(I_{\mathcal{U}_j}+\lambda_j^{-1}N_j\right)$ with $\lambda_j^{-1}N_j$ nilpotent. By Lemma \ref{lem:identity-plus-nilpotent-square-root}, choose $B_j^2=I_{\mathcal{U}_j}+\lambda_j^{-1}N_j$; choose $\alpha_j^2=\lambda_j$ and set $R_j=\alpha_jB_j$. Then $R_j^2=T|_{\mathcal{U}_j}$. Define $R=R_1\oplus\cdots\oplus R_m$ on the direct sum. Then $R^2=T$.\\
\hlt{Structural Commentary}: Take square roots on primary blocks, then recombine.\\
\hlt{Key Insight}: Invertibility makes every scalar block nonzero.
\end{proof}
\end{theorem}
